
\documentclass{article}
\usepackage{colortbl}
\usepackage{makecell}
\usepackage{multirow}
\usepackage{supertabular}

\begin{document}

\newcounter{utterance}

\centering \large Interaction Transcript for game `hot\_air\_balloon', experiment `air\_balloon\_survival\_en\_complexity\_hard', episode 5 with mistral{-}3{-}large{-}2512{-}t1.0.
\vspace{24pt}

{ \footnotesize  \setcounter{utterance}{1}
\setlength{\tabcolsep}{0pt}
\begin{supertabular}{c@{$\;$}|p{.15\linewidth}@{}p{.15\linewidth}p{.15\linewidth}p{.15\linewidth}p{.15\linewidth}p{.15\linewidth}}
    \# & $\;$A & \multicolumn{4}{c}{Game Master} & $\;\:$B\\
    \hline

    \theutterance \stepcounter{utterance}  
    & & \multicolumn{4}{p{0.6\linewidth}}{
        \cellcolor[rgb]{0.9,0.9,0.9}{
            \makecell[{{p{\linewidth}}}]{
                \texttt{\tiny{[P1$\langle$GM]}}
                \texttt{You are participating in a collaborative negotiation game.} \\
\\ 
\texttt{Together with another participant, you must agree on a single set of items that will be kept. Each of you has your own view of how much each item matters to you (importance). You do not know how the other participant values the items. Additionally, you are given the effort each item demands.} \\
\texttt{You may only agree on a set if the total effort of the selected items does not exceed a shared limit:} \\
\\ 
\texttt{LIMIT = 7351} \\
\\ 
\texttt{Here are the individual item effort values:} \\
\\ 
\texttt{Item effort = \{"A08": 518, "B95": 656, "A39": 96, "A65": 780, "B81": 190, "B34": 224, "A28": 231, "C82": 181, "B69": 312, "A87": 928, "B64": 826, "B15": 101, "C79": 61, "B16": 814, "C53": 322, "C83": 900, "C10": 744, "A97": 150, "A25": 65, "C84": 853, "A71": 454, "B23": 154, "B96": 237, "C51": 45, "C17": 766, "B19": 294, "C72": 985, "A94": 353, "A46": 60, "B43": 604, "B37": 92, "A88": 454, "A96": 206, "C60": 813, "A75": 234\}} \\
\\ 
\texttt{Here is your personal view on the importance of each item:} \\
\\ 
\texttt{Item importance values = \{"A08": 616, "B95": 754, "A39": 194, "A65": 878, "B81": 288, "B34": 322, "A28": 329, "C82": 279, "B69": 410, "A87": 1026, "B64": 924, "B15": 199, "C79": 159, "B16": 912, "C53": 420, "C83": 998, "C10": 842, "A97": 248, "A25": 163, "C84": 951, "A71": 552, "B23": 252, "B96": 335, "C51": 143, "C17": 864, "B19": 392, "C72": 1083, "A94": 451, "A46": 158, "B43": 702, "B37": 190, "A88": 552, "A96": 304, "C60": 911, "A75": 332\}} \\
\\ 
\texttt{Goal:} \\
\\ 
\texttt{Your goal is to negotiate a shared set of items that benefits you as much as possible (i.e., maximizes total importance to YOU), while staying within the LIMIT. You are not required to make a PROPOSAL in every message {-} you can simply negotiate as well. All tactics are allowed!} \\
\\ 
\texttt{Interaction Protocol:} \\
\\ 
\texttt{You may only use the following structured formats in a message:} \\
\\ 
\texttt{PROPOSAL: \{'A', 'B', 'C', …\}} \\
\texttt{Propose keeping exactly those items.} \\
\texttt{REFUSE: \{'A', 'B', 'C', …\}} \\
\texttt{Explicitly reject opponent's proposal.} \\
\texttt{ARGUMENT: \{'...'\}} \\
\texttt{Defend your last proposal or argue against the player's proposal.} \\
\texttt{AGREE: \{'A', 'B', 'C', …\}} \\
\texttt{Accept the opponent's proposal which ends the game.} \\
\texttt{STRATEGIC REASONING: \{'...'\}} \\
\texttt{	Describe your strategic reasoning or anticipation explaining your choice of action. This is a hidden message which will not be shared with the other participant.} \\
\\ 
\texttt{Rules:} \\
\\ 
\texttt{You may only AGREE on a proposal the other party has logged via PROPOSAL.} \\
\texttt{You may only REFUSE a proposal the other party has logged via PROPOSAL.} \\
\texttt{Total effort of any PROPOSAL or AGREE set must be ≤ LIMIT.} \\
\texttt{Do NOT reveal your hidden importance scores.} \\
\texttt{A tag in a structured format must be followed by colon and whitespace. The argument must be a python set containing 0 or more strings.} \\
\texttt{So, it must be of the form TAG: \{...\}} \\
\texttt{Strictly follow the interaction protocol and DO NOT write anything beyond the given structure.} \\
\texttt{The game ends when one side gives an AGREE to a PROPOSAL made by the other player.} \\
\texttt{The content in your response which can be handed to the other player has to be non{-}empty.} \\
\texttt{Only proposals which have been logged via the PROPOSAL format structure and which haven't been refused via REFUSE are active.} \\
\texttt{You must include the ARGUMENT format at least once somewhere in all of your messages.} \\
\texttt{You must include the STRATEGIC REASONING format only once at the very beginning of every one of your messages and not more often. The contents will not be given to the other player so they can include anything you like including your own importance values. Here you should reason step by step to come up with you next move.} \\
            }
        }
    }
    & \\ \\

    \theutterance \stepcounter{utterance}  
    & \multicolumn{4}{p{0.6\linewidth}}{
        \cellcolor[rgb]{0.9,0.9,0.9}{
            \makecell[{{p{\linewidth}}}]{
                \texttt{\tiny{[P1$\rangle$GM]}}
                \texttt{STRATEGIC REASONING: \{} \\
\texttt{I will start by analyzing the items to identify the most valuable ones for me (high importance, low effort) and the least valuable ones (low importance, high effort). This will help me prioritize items in negotiations and identify potential trade{-}offs.} \\
\\ 
\texttt{Key observations:} \\
\texttt{1. High importance, low effort items (top priorities):} \\
\texttt{   {-} A46 (60 effort, 158 importance)} \\
\texttt{   {-} C79 (61 effort, 159 importance)} \\
\texttt{   {-} A25 (65 effort, 163 importance)} \\
\texttt{   {-} A39 (96 effort, 194 importance)} \\
\texttt{   {-} B15 (101 effort, 199 importance)} \\
\texttt{   {-} B37 (92 effort, 190 importance) {-} slightly worse ratio but still good} \\
\texttt{   {-} A97 (150 effort, 248 importance)} \\
\texttt{   {-} B81 (190 effort, 288 importance)} \\
\texttt{   {-} B23 (154 effort, 252 importance)} \\
\\ 
\texttt{2. Low importance, high effort items (avoid):} \\
\texttt{   {-} A87 (928 effort, 1026 importance) {-} high importance but very high effort} \\
\texttt{   {-} C72 (985 effort, 1083 importance) {-} same as above} \\
\texttt{   {-} C83 (900 effort, 998 importance)} \\
\texttt{   {-} C84 (853 effort, 951 importance)} \\
\texttt{   {-} B64 (826 effort, 924 importance)} \\
\texttt{   {-} C60 (813 effort, 911 importance)} \\
\texttt{   {-} B16 (814 effort, 912 importance)} \\
\texttt{   {-} C10 (744 effort, 842 importance)} \\
\texttt{   {-} C17 (766 effort, 864 importance)} \\
\\ 
\texttt{3. Strategy:} \\
\texttt{   {-} I should propose a set that includes as many high{-}importance, low{-}effort items as possible while staying under the limit.} \\
\texttt{   {-} I should avoid including the very high{-}effort items unless they are absolutely necessary for negotiation.} \\
\texttt{   {-} I will start with a strong proposal that maximizes my importance while leaving room for negotiation.} \\
\texttt{   {-} I will argue that certain items are essential due to their low effort and high value} \\
            }
        }
    }
    & & \\ \\

    \theutterance \stepcounter{utterance}  
    & & \multicolumn{4}{p{0.6\linewidth}}{
        \cellcolor[rgb]{0.9,0.9,0.9}{
            \makecell[{{p{\linewidth}}}]{
                \texttt{\tiny{[P1$\langle$GM]}}
                \texttt{Your response did not start with the proper strategic reasoning tag at the very beginning of your response. The very first structured format must be of the form STRATEGIC REASONING: \{...\}. Try again.} \\
            }
        }
    }
    & \\ \\

    \theutterance \stepcounter{utterance}  
    & \multicolumn{4}{p{0.6\linewidth}}{
        \cellcolor[rgb]{0.9,0.9,0.9}{
            \makecell[{{p{\linewidth}}}]{
                \texttt{\tiny{[P1$\rangle$GM]}}
                \texttt{STRATEGIC REASONING: \{} \\
\texttt{Initial analysis:} \\
\texttt{1. My goal is to maximize my total importance while staying under the LIMIT (7351 effort).} \\
\texttt{2. I will prioritize items with the best importance{-}to{-}effort ratio (importance/effort).} \\
\texttt{3. Top 10 items by ratio (importance/effort):} \\
\texttt{   {-} A46: 2.63 (60/158)} \\
\texttt{   {-} C51: 3.18 (45/143)} \\
\texttt{   {-} C79: 2.61 (61/159)} \\
\texttt{   {-} A25: 2.51 (65/163)} \\
\texttt{   {-} A39: 2.02 (96/194)} \\
\texttt{   {-} B37: 2.07 (92/190)} \\
\texttt{   {-} B15: 1.97 (101/199)} \\
\texttt{   {-} A97: 1.65 (150/248)} \\
\texttt{   {-} B81: 1.52 (190/288)} \\
\texttt{   {-} B23: 1.64 (154/252)} \\
\\ 
\texttt{4. Proposal strategy:} \\
\texttt{   {-} Start with a proposal that includes all top{-}ratio items plus some high{-}importance items.} \\
\texttt{   {-} Exclude the worst{-}ratio items (especially those with effort > 800).} \\
\texttt{   {-} Calculate total effort to ensure it's under LIMIT while maximizing importance.} \\
\texttt{   {-} Leave room for negotiation by including some mid{-}value items.} \\
\\ 
\texttt{5. Initial proposal calculation:} \\
\texttt{   {-} Core items (top ratios): A46, C51, C79, A25, A39, B37, B15, A97, B81, B23 (total effort: 915)} \\
\texttt{   {-} Add high{-}importance items: B95 (656/754), A65 (780/878), C53 (322/420), B69 (312/410), A28 (231/329), B34 (224/322) (additional effort: 2525)} \\
\texttt{   {-} Total effort so far:} \\
            }
        }
    }
    & & \\ \\

    \theutterance \stepcounter{utterance}  
    & & \multicolumn{4}{p{0.6\linewidth}}{
        \cellcolor[rgb]{0.9,0.9,0.9}{
            \makecell[{{p{\linewidth}}}]{
                \texttt{\tiny{[P1$\langle$GM]}}
                \texttt{Your response did not start with the proper strategic reasoning tag at the very beginning of your response. The very first structured format must be of the form STRATEGIC REASONING: \{...\}. Try again.} \\
            }
        }
    }
    & \\ \\

    \theutterance \stepcounter{utterance}  
    & \multicolumn{4}{p{0.6\linewidth}}{
        \cellcolor[rgb]{0.9,0.9,0.9}{
            \makecell[{{p{\linewidth}}}]{
                \texttt{\tiny{[P1$\rangle$GM]}}
                \texttt{STRATEGIC REASONING: \{} \\
\texttt{1. **Objective**: Maximize my total importance while keeping total effort ≤ 7351.} \\
\texttt{2. **Key Insight**: Prioritize items with the highest importance{-}to{-}effort ratio (importance/effort) to maximize value per unit of effort.} \\
\texttt{3. **Top{-}Ratio Items** (importance/effort > 1.5):} \\
\texttt{   {-} A46 (2.63), C51 (3.18), C79 (2.61), A25 (2.51), A39 (2.02), B37 (2.07), B15 (1.97), A97 (1.65), B81 (1.52), B23 (1.64)} \\
\texttt{   {-} These are non{-}negotiable for me due to their efficiency.} \\
\texttt{4. **High{-}Importance Items** (importance > 800, but high effort):} \\
\texttt{   {-} A65 (878), B95 (754), C83 (998), C84 (951), B64 (924), C60 (911), B16 (912)} \\
\texttt{   {-} These are valuable but costly; I’ll only include them if the opponent concedes other high{-}value items.} \\
\texttt{5. **Proposal Strategy**:} \\
\texttt{   {-} Start with a strong proposal including all top{-}ratio items + select high{-}importance items, leaving \textasciitilde{}1000 effort buffer for negotiation.} \\
\texttt{   {-} Exclude the worst{-}ratio items (e.g., C72, A87, C83, C84) to avoid wasting effort.} \\
\texttt{   {-} Total effort of top{-}ratio items: 915 (A46, C51, C79, A25, A39, B37, B15, A97, B81, B23).} \\
\texttt{   {-} Add high{-}importance items: B95 (656), A65 (780), C53 (322), B69 (312), A28 (231), B34 (224) → Total effort = 3440.} \\
\texttt{   {-} Remaining effort: 7351 {-}} \\
            }
        }
    }
    & & \\ \\

    \theutterance \stepcounter{utterance}  
    & & & \multicolumn{2}{p{0.3\linewidth}}{
        \cellcolor[rgb]{0.9,0.9,0.9}{
            \makecell[{{p{\linewidth}}}]{
                \texttt{\tiny{[GM$|$GM]}}
                \texttt{abort} \\
            }
        }
    }
    & & \\ \\

    \theutterance \stepcounter{utterance}  
    & & & \multicolumn{2}{p{0.3\linewidth}}{
        \cellcolor[rgb]{0.9,0.9,0.9}{
            \makecell[{{p{\linewidth}}}]{
                \texttt{\tiny{[GM$|$GM]}}
                \texttt{end game} \\
            }
        }
    }
    & & \\ \\

\end{supertabular}
}

\end{document}
